\documentclass[border=10pt,tikz]{standalone}
\usepackage{fontspec}
\setmainfont{Apple SD Gothic Neo}
\setsansfont{Apple SD Gothic Neo}
\usepackage{tikz}
\usetikzlibrary{positioning,arrows.meta,shapes,fit,backgrounds,calc}
\begin{document}
% _design.tex — shared TikZ design system for all blog diagrams
% Inputs nothing on its own — meant to be \input{} from each diagram .tex
% AFTER \begin{document}, BEFORE \begin{tikzpicture}.
%
% Companion files:
%   _design-math.tex — math diagrams (PGFPlots, book-notes palette)
%
% Provides a unified visual language across GoF / DSA / EMC++ / Beautiful / 시리즈:
%   - Colors: muted, accessible, color-coded by role
%   - Node styles: consistent sizing, rounded, subtle borders
%   - Edge styles: UML-style inheritance, aggregation, plus dashed alternates
%   - Annotation: callouts, notes, badges
%
% Usage:
%   \documentclass[border=10pt,tikz]{standalone}
%   \usepackage{tikz}
%   \usetikzlibrary{positioning,arrows.meta,shapes,fit,backgrounds,calc,decorations.pathreplacing}
%   \begin{document}
%   % _design.tex — shared TikZ design system for all blog diagrams
% Inputs nothing on its own — meant to be \input{} from each diagram .tex
% AFTER \begin{document}, BEFORE \begin{tikzpicture}.
%
% Companion files:
%   _design-math.tex — math diagrams (PGFPlots, book-notes palette)
%
% Provides a unified visual language across GoF / DSA / EMC++ / Beautiful / 시리즈:
%   - Colors: muted, accessible, color-coded by role
%   - Node styles: consistent sizing, rounded, subtle borders
%   - Edge styles: UML-style inheritance, aggregation, plus dashed alternates
%   - Annotation: callouts, notes, badges
%
% Usage:
%   \documentclass[border=10pt,tikz]{standalone}
%   \usepackage{tikz}
%   \usetikzlibrary{positioning,arrows.meta,shapes,fit,backgrounds,calc,decorations.pathreplacing}
%   \begin{document}
%   % _design.tex — shared TikZ design system for all blog diagrams
% Inputs nothing on its own — meant to be \input{} from each diagram .tex
% AFTER \begin{document}, BEFORE \begin{tikzpicture}.
%
% Companion files:
%   _design-math.tex — math diagrams (PGFPlots, book-notes palette)
%
% Provides a unified visual language across GoF / DSA / EMC++ / Beautiful / 시리즈:
%   - Colors: muted, accessible, color-coded by role
%   - Node styles: consistent sizing, rounded, subtle borders
%   - Edge styles: UML-style inheritance, aggregation, plus dashed alternates
%   - Annotation: callouts, notes, badges
%
% Usage:
%   \documentclass[border=10pt,tikz]{standalone}
%   \usepackage{tikz}
%   \usetikzlibrary{positioning,arrows.meta,shapes,fit,backgrounds,calc,decorations.pathreplacing}
%   \begin{document}
%   \input{../_design.tex}     % adjust relative path
%   \begin{tikzpicture}[blog]  % apply 'blog' style
%   ...
%
% Backward-compat note: This file only ADDS to the public API. Existing styles
% (abs/con/imp/cli/hi/vx/q/inh/agg/uses/arr/edge/alt/ret/thr/group/glabel) are
% preserved. New helpers are documented below.

% ─── Core palette (fills paired with darker borders) ─────────
\definecolor{absfill}{RGB}{220,232,247}   % cool blue
\definecolor{absbord}{RGB}{72,118,182}
\definecolor{confill}{RGB}{249,222,222}   % warm red
\definecolor{conbord}{RGB}{197,93,93}
\definecolor{impfill}{RGB}{223,238,222}   % soft green
\definecolor{impbord}{RGB}{99,154,93}
\definecolor{clifill}{RGB}{251,239,205}   % pale yellow
\definecolor{clibord}{RGB}{200,167,72}
\definecolor{hifill} {RGB}{253,224,200}   % accent orange
\definecolor{hibord} {RGB}{210,135,68}
\definecolor{nodefill}{RGB}{250,250,250}
\definecolor{nodebord}{RGB}{120,120,120}
\definecolor{edgec}  {RGB}{80,80,80}
\definecolor{edgesc} {RGB}{145,145,145}
\definecolor{groupbord}{RGB}{200,200,200}

% ─── Extended palette (new) ──────────────────────────────────
\definecolor{prpfill}{RGB}{233,224,247}   % muted purple (state / interface)
\definecolor{prpbord}{RGB}{132,103,182}
\definecolor{tealfill}{RGB}{215,235,235}  % teal (services / utility)
\definecolor{tealbord}{RGB}{ 78,135,140}
\definecolor{warnfill}{RGB}{251,229,229}  % deeper warn (error states)
\definecolor{warnbord}{RGB}{180, 55, 55}
\definecolor{okfill}  {RGB}{220,240,225}  % success
\definecolor{okbord}  {RGB}{ 75,140, 90}
\definecolor{neutfill}{RGB}{238,238,238}  % neutral panel
\definecolor{neutbord}{RGB}{170,170,170}
\definecolor{inkc}    {RGB}{ 30, 30, 30}  % strong ink (titles)
\definecolor{mutec}   {RGB}{120,120,120}  % muted text
\definecolor{focusc}  {RGB}{210, 70, 70}  % focus / highlight accent
\definecolor{gridc}   {RGB}{220,220,220}  % background grid

\tikzset{
  blog/.style={
    font=\sffamily\small\linespread{1.18}\selectfont,
    every node/.style={font=\sffamily\small\linespread{1.18}\selectfont}
  },
  % ─── existing node roles (unchanged) ───
  cls/.style={
    rectangle, rounded corners=2pt,
    draw=nodebord, line width=0.4pt,
    minimum height=8mm, minimum width=28mm,
    inner sep=3pt, fill=nodefill, align=center
  },
  abs/.style={cls, draw=absbord, fill=absfill, font=\sffamily\itshape\small},
  con/.style={cls, draw=conbord, fill=confill},
  imp/.style={cls, draw=impbord, fill=impfill},
  cli/.style={cls, draw=clibord, fill=clifill},
  hi/.style={cls, draw=hibord, fill=hifill, font=\sffamily\bfseries\small},
  % ─── new node roles ───
  prp/.style={cls, draw=prpbord, fill=prpfill},               % protocol / interface / state
  tealn/.style={cls, draw=tealbord, fill=tealfill},           % service / utility
  warn/.style={cls, draw=warnbord, fill=warnfill, font=\sffamily\bfseries\small},  % error state
  ok/.style={cls, draw=okbord, fill=okfill},                  % success state
  neut/.style={cls, draw=neutbord, fill=neutfill},            % neutral
  % ─── size variants ───
  clsSm/.style={cls, minimum height=6mm, minimum width=20mm, font=\sffamily\footnotesize, inner sep=2pt},
  clsLg/.style={cls, minimum height=11mm, minimum width=36mm, font=\sffamily\normalsize},
  absSm/.style={clsSm, draw=absbord, fill=absfill, font=\sffamily\itshape\footnotesize},
  conSm/.style={clsSm, draw=conbord, fill=confill},
  impSm/.style={clsSm, draw=impbord, fill=impfill},
  % ─── circle (existing) ───
  vx/.style={
    circle, draw=absbord, line width=0.4pt,
    minimum size=8mm, fill=absfill, inner sep=0pt
  },
  vxc/.style={vx, draw=conbord, fill=confill},
  vxh/.style={vx, draw=hibord, fill=hifill, font=\sffamily\bfseries},
  % new circle variants
  vxi/.style={vx, draw=impbord, fill=impfill},                % impl green
  vxp/.style={vx, draw=prpbord, fill=prpfill},                % protocol purple
  vxw/.style={vx, draw=warnbord, fill=warnfill},              % warn red
  vxn/.style={vx, draw=nodebord, fill=nodefill},              % neutral
  vxSm/.style={vx, minimum size=6mm, font=\sffamily\footnotesize},
  vxLg/.style={vx, minimum size=11mm, font=\sffamily\normalsize},
  % ─── decision diamond (existing) ───
  q/.style={
    diamond, draw=clibord, line width=0.4pt, fill=clifill,
    inner sep=2pt, aspect=2.4, align=center
  },
  % rounded callout / note (new)
  note/.style={
    rectangle, rounded corners=3pt,
    draw=clibord, line width=0.4pt, fill=clifill,
    inner sep=3pt, font=\sffamily\footnotesize, align=left
  },
  badge/.style={                                              % small inline marker
    rectangle, rounded corners=2pt,
    draw=hibord, line width=0.3pt, fill=hifill,
    inner sep=1.5pt, font=\sffamily\scriptsize\bfseries
  },
  % ─── edges (existing) ───
  inh/.style={                                  % UML inheritance
    -{Triangle[length=3mm,width=3mm,open]},
    draw=edgec, line width=0.5pt
  },
  agg/.style={                                  % UML aggregation
    {Diamond[length=2.5mm,width=2.5mm,open]}-{Stealth[length=2.2mm]},
    draw=edgec, line width=0.45pt
  },
  uses/.style={                                 % directed reference
    -{Stealth[length=2.2mm]}, draw=edgec, line width=0.5pt
  },
  arr/.style={                                  % plain arrow (graphs/trees)
    -{Stealth[length=2mm]}, draw=edgec, line width=0.45pt
  },
  edge/.style={                                 % undirected (trees)
    draw=edgec, line width=0.45pt
  },
  alt/.style={                                  % dashed: similar/alt
    -{Stealth[length=2mm]}, dashed, draw=edgesc, line width=0.45pt
  },
  ret/.style={                                  % dashed creates/returns
    -{Stealth[length=2mm]}, dashed, draw=edgesc, line width=0.45pt
  },
  thr/.style={                                  % threading link (red dashed)
    ->, dashed, draw=conbord, line width=0.5pt
  },
  % ─── new edges ───
  comp/.style={                                 % UML composition (filled diamond)
    {Diamond[length=2.5mm,width=2.5mm]}-{Stealth[length=2.2mm]},
    draw=edgec, line width=0.45pt
  },
  realize/.style={                              % UML realize (impl interface)
    -{Triangle[length=3mm,width=3mm,open]}, dashed,
    draw=edgec, line width=0.5pt
  },
  msgsync/.style={                              % synchronous message (sequence)
    -{Stealth[length=2.6mm]}, draw=edgec, line width=0.55pt
  },
  msgasync/.style={                             % asynchronous message
    -{Stealth[length=2.4mm,inset=0.4mm,fill=edgec]},
    draw=edgec, line width=0.5pt
  },
  msgret/.style={                               % return message (dashed)
    -{Stealth[length=2.2mm]}, dashed, draw=edgesc, line width=0.45pt
  },
  flow/.style={                                 % control flow (thicker, hi-vis)
    -{Stealth[length=2.4mm]}, draw=inkc, line width=0.7pt
  },
  focus/.style={                                % focus / hot path
    -{Stealth[length=2.4mm]}, draw=focusc, line width=0.8pt
  },
  selfloop/.style={                             % self loop on a node
    ->, draw=edgec, line width=0.45pt, in=120, out=60, loop, looseness=8
  },
  % ─── group / region (existing + new) ───
  group/.style={
    rectangle, rounded corners=4pt,
    draw=groupbord, line width=0.4pt,
    inner sep=10pt, fill=black!2
  },
  glabel/.style={font=\sffamily\bfseries\small, anchor=north west, inner sep=2pt},
  panel/.style={                                % stronger panel for layered diagrams
    rectangle, rounded corners=4pt,
    draw=neutbord, line width=0.45pt,
    inner sep=10pt, fill=neutfill!50
  },
  panelTitle/.style={font=\sffamily\bfseries\small, text=inkc, anchor=north west, inner sep=2pt},
  panelSub/.style={font=\sffamily\footnotesize, text=mutec, anchor=north west, inner sep=2pt},
  % ─── typography ───
  caption/.style={font=\sffamily\footnotesize\itshape, text=mutec, align=center},
  legendkey/.style={rectangle, minimum width=10mm, minimum height=3mm, inner sep=0pt, line width=0.3pt},
  % ─── grid / coordinate hints ───
  bggrid/.style={draw=gridc, line width=0.3pt, dotted},
}

% Korean font convenience macro — included by every Hangul diagram
\providecommand{\KOR}{}
     % adjust relative path
%   \begin{tikzpicture}[blog]  % apply 'blog' style
%   ...
%
% Backward-compat note: This file only ADDS to the public API. Existing styles
% (abs/con/imp/cli/hi/vx/q/inh/agg/uses/arr/edge/alt/ret/thr/group/glabel) are
% preserved. New helpers are documented below.

% ─── Core palette (fills paired with darker borders) ─────────
\definecolor{absfill}{RGB}{220,232,247}   % cool blue
\definecolor{absbord}{RGB}{72,118,182}
\definecolor{confill}{RGB}{249,222,222}   % warm red
\definecolor{conbord}{RGB}{197,93,93}
\definecolor{impfill}{RGB}{223,238,222}   % soft green
\definecolor{impbord}{RGB}{99,154,93}
\definecolor{clifill}{RGB}{251,239,205}   % pale yellow
\definecolor{clibord}{RGB}{200,167,72}
\definecolor{hifill} {RGB}{253,224,200}   % accent orange
\definecolor{hibord} {RGB}{210,135,68}
\definecolor{nodefill}{RGB}{250,250,250}
\definecolor{nodebord}{RGB}{120,120,120}
\definecolor{edgec}  {RGB}{80,80,80}
\definecolor{edgesc} {RGB}{145,145,145}
\definecolor{groupbord}{RGB}{200,200,200}

% ─── Extended palette (new) ──────────────────────────────────
\definecolor{prpfill}{RGB}{233,224,247}   % muted purple (state / interface)
\definecolor{prpbord}{RGB}{132,103,182}
\definecolor{tealfill}{RGB}{215,235,235}  % teal (services / utility)
\definecolor{tealbord}{RGB}{ 78,135,140}
\definecolor{warnfill}{RGB}{251,229,229}  % deeper warn (error states)
\definecolor{warnbord}{RGB}{180, 55, 55}
\definecolor{okfill}  {RGB}{220,240,225}  % success
\definecolor{okbord}  {RGB}{ 75,140, 90}
\definecolor{neutfill}{RGB}{238,238,238}  % neutral panel
\definecolor{neutbord}{RGB}{170,170,170}
\definecolor{inkc}    {RGB}{ 30, 30, 30}  % strong ink (titles)
\definecolor{mutec}   {RGB}{120,120,120}  % muted text
\definecolor{focusc}  {RGB}{210, 70, 70}  % focus / highlight accent
\definecolor{gridc}   {RGB}{220,220,220}  % background grid

\tikzset{
  blog/.style={
    font=\sffamily\small\linespread{1.18}\selectfont,
    every node/.style={font=\sffamily\small\linespread{1.18}\selectfont}
  },
  % ─── existing node roles (unchanged) ───
  cls/.style={
    rectangle, rounded corners=2pt,
    draw=nodebord, line width=0.4pt,
    minimum height=8mm, minimum width=28mm,
    inner sep=3pt, fill=nodefill, align=center
  },
  abs/.style={cls, draw=absbord, fill=absfill, font=\sffamily\itshape\small},
  con/.style={cls, draw=conbord, fill=confill},
  imp/.style={cls, draw=impbord, fill=impfill},
  cli/.style={cls, draw=clibord, fill=clifill},
  hi/.style={cls, draw=hibord, fill=hifill, font=\sffamily\bfseries\small},
  % ─── new node roles ───
  prp/.style={cls, draw=prpbord, fill=prpfill},               % protocol / interface / state
  tealn/.style={cls, draw=tealbord, fill=tealfill},           % service / utility
  warn/.style={cls, draw=warnbord, fill=warnfill, font=\sffamily\bfseries\small},  % error state
  ok/.style={cls, draw=okbord, fill=okfill},                  % success state
  neut/.style={cls, draw=neutbord, fill=neutfill},            % neutral
  % ─── size variants ───
  clsSm/.style={cls, minimum height=6mm, minimum width=20mm, font=\sffamily\footnotesize, inner sep=2pt},
  clsLg/.style={cls, minimum height=11mm, minimum width=36mm, font=\sffamily\normalsize},
  absSm/.style={clsSm, draw=absbord, fill=absfill, font=\sffamily\itshape\footnotesize},
  conSm/.style={clsSm, draw=conbord, fill=confill},
  impSm/.style={clsSm, draw=impbord, fill=impfill},
  % ─── circle (existing) ───
  vx/.style={
    circle, draw=absbord, line width=0.4pt,
    minimum size=8mm, fill=absfill, inner sep=0pt
  },
  vxc/.style={vx, draw=conbord, fill=confill},
  vxh/.style={vx, draw=hibord, fill=hifill, font=\sffamily\bfseries},
  % new circle variants
  vxi/.style={vx, draw=impbord, fill=impfill},                % impl green
  vxp/.style={vx, draw=prpbord, fill=prpfill},                % protocol purple
  vxw/.style={vx, draw=warnbord, fill=warnfill},              % warn red
  vxn/.style={vx, draw=nodebord, fill=nodefill},              % neutral
  vxSm/.style={vx, minimum size=6mm, font=\sffamily\footnotesize},
  vxLg/.style={vx, minimum size=11mm, font=\sffamily\normalsize},
  % ─── decision diamond (existing) ───
  q/.style={
    diamond, draw=clibord, line width=0.4pt, fill=clifill,
    inner sep=2pt, aspect=2.4, align=center
  },
  % rounded callout / note (new)
  note/.style={
    rectangle, rounded corners=3pt,
    draw=clibord, line width=0.4pt, fill=clifill,
    inner sep=3pt, font=\sffamily\footnotesize, align=left
  },
  badge/.style={                                              % small inline marker
    rectangle, rounded corners=2pt,
    draw=hibord, line width=0.3pt, fill=hifill,
    inner sep=1.5pt, font=\sffamily\scriptsize\bfseries
  },
  % ─── edges (existing) ───
  inh/.style={                                  % UML inheritance
    -{Triangle[length=3mm,width=3mm,open]},
    draw=edgec, line width=0.5pt
  },
  agg/.style={                                  % UML aggregation
    {Diamond[length=2.5mm,width=2.5mm,open]}-{Stealth[length=2.2mm]},
    draw=edgec, line width=0.45pt
  },
  uses/.style={                                 % directed reference
    -{Stealth[length=2.2mm]}, draw=edgec, line width=0.5pt
  },
  arr/.style={                                  % plain arrow (graphs/trees)
    -{Stealth[length=2mm]}, draw=edgec, line width=0.45pt
  },
  edge/.style={                                 % undirected (trees)
    draw=edgec, line width=0.45pt
  },
  alt/.style={                                  % dashed: similar/alt
    -{Stealth[length=2mm]}, dashed, draw=edgesc, line width=0.45pt
  },
  ret/.style={                                  % dashed creates/returns
    -{Stealth[length=2mm]}, dashed, draw=edgesc, line width=0.45pt
  },
  thr/.style={                                  % threading link (red dashed)
    ->, dashed, draw=conbord, line width=0.5pt
  },
  % ─── new edges ───
  comp/.style={                                 % UML composition (filled diamond)
    {Diamond[length=2.5mm,width=2.5mm]}-{Stealth[length=2.2mm]},
    draw=edgec, line width=0.45pt
  },
  realize/.style={                              % UML realize (impl interface)
    -{Triangle[length=3mm,width=3mm,open]}, dashed,
    draw=edgec, line width=0.5pt
  },
  msgsync/.style={                              % synchronous message (sequence)
    -{Stealth[length=2.6mm]}, draw=edgec, line width=0.55pt
  },
  msgasync/.style={                             % asynchronous message
    -{Stealth[length=2.4mm,inset=0.4mm,fill=edgec]},
    draw=edgec, line width=0.5pt
  },
  msgret/.style={                               % return message (dashed)
    -{Stealth[length=2.2mm]}, dashed, draw=edgesc, line width=0.45pt
  },
  flow/.style={                                 % control flow (thicker, hi-vis)
    -{Stealth[length=2.4mm]}, draw=inkc, line width=0.7pt
  },
  focus/.style={                                % focus / hot path
    -{Stealth[length=2.4mm]}, draw=focusc, line width=0.8pt
  },
  selfloop/.style={                             % self loop on a node
    ->, draw=edgec, line width=0.45pt, in=120, out=60, loop, looseness=8
  },
  % ─── group / region (existing + new) ───
  group/.style={
    rectangle, rounded corners=4pt,
    draw=groupbord, line width=0.4pt,
    inner sep=10pt, fill=black!2
  },
  glabel/.style={font=\sffamily\bfseries\small, anchor=north west, inner sep=2pt},
  panel/.style={                                % stronger panel for layered diagrams
    rectangle, rounded corners=4pt,
    draw=neutbord, line width=0.45pt,
    inner sep=10pt, fill=neutfill!50
  },
  panelTitle/.style={font=\sffamily\bfseries\small, text=inkc, anchor=north west, inner sep=2pt},
  panelSub/.style={font=\sffamily\footnotesize, text=mutec, anchor=north west, inner sep=2pt},
  % ─── typography ───
  caption/.style={font=\sffamily\footnotesize\itshape, text=mutec, align=center},
  legendkey/.style={rectangle, minimum width=10mm, minimum height=3mm, inner sep=0pt, line width=0.3pt},
  % ─── grid / coordinate hints ───
  bggrid/.style={draw=gridc, line width=0.3pt, dotted},
}

% Korean font convenience macro — included by every Hangul diagram
\providecommand{\KOR}{}
     % adjust relative path
%   \begin{tikzpicture}[blog]  % apply 'blog' style
%   ...
%
% Backward-compat note: This file only ADDS to the public API. Existing styles
% (abs/con/imp/cli/hi/vx/q/inh/agg/uses/arr/edge/alt/ret/thr/group/glabel) are
% preserved. New helpers are documented below.

% ─── Core palette (fills paired with darker borders) ─────────
\definecolor{absfill}{RGB}{220,232,247}   % cool blue
\definecolor{absbord}{RGB}{72,118,182}
\definecolor{confill}{RGB}{249,222,222}   % warm red
\definecolor{conbord}{RGB}{197,93,93}
\definecolor{impfill}{RGB}{223,238,222}   % soft green
\definecolor{impbord}{RGB}{99,154,93}
\definecolor{clifill}{RGB}{251,239,205}   % pale yellow
\definecolor{clibord}{RGB}{200,167,72}
\definecolor{hifill} {RGB}{253,224,200}   % accent orange
\definecolor{hibord} {RGB}{210,135,68}
\definecolor{nodefill}{RGB}{250,250,250}
\definecolor{nodebord}{RGB}{120,120,120}
\definecolor{edgec}  {RGB}{80,80,80}
\definecolor{edgesc} {RGB}{145,145,145}
\definecolor{groupbord}{RGB}{200,200,200}

% ─── Extended palette (new) ──────────────────────────────────
\definecolor{prpfill}{RGB}{233,224,247}   % muted purple (state / interface)
\definecolor{prpbord}{RGB}{132,103,182}
\definecolor{tealfill}{RGB}{215,235,235}  % teal (services / utility)
\definecolor{tealbord}{RGB}{ 78,135,140}
\definecolor{warnfill}{RGB}{251,229,229}  % deeper warn (error states)
\definecolor{warnbord}{RGB}{180, 55, 55}
\definecolor{okfill}  {RGB}{220,240,225}  % success
\definecolor{okbord}  {RGB}{ 75,140, 90}
\definecolor{neutfill}{RGB}{238,238,238}  % neutral panel
\definecolor{neutbord}{RGB}{170,170,170}
\definecolor{inkc}    {RGB}{ 30, 30, 30}  % strong ink (titles)
\definecolor{mutec}   {RGB}{120,120,120}  % muted text
\definecolor{focusc}  {RGB}{210, 70, 70}  % focus / highlight accent
\definecolor{gridc}   {RGB}{220,220,220}  % background grid

\tikzset{
  blog/.style={
    font=\sffamily\small\linespread{1.18}\selectfont,
    every node/.style={font=\sffamily\small\linespread{1.18}\selectfont}
  },
  % ─── existing node roles (unchanged) ───
  cls/.style={
    rectangle, rounded corners=2pt,
    draw=nodebord, line width=0.4pt,
    minimum height=8mm, minimum width=28mm,
    inner sep=3pt, fill=nodefill, align=center
  },
  abs/.style={cls, draw=absbord, fill=absfill, font=\sffamily\itshape\small},
  con/.style={cls, draw=conbord, fill=confill},
  imp/.style={cls, draw=impbord, fill=impfill},
  cli/.style={cls, draw=clibord, fill=clifill},
  hi/.style={cls, draw=hibord, fill=hifill, font=\sffamily\bfseries\small},
  % ─── new node roles ───
  prp/.style={cls, draw=prpbord, fill=prpfill},               % protocol / interface / state
  tealn/.style={cls, draw=tealbord, fill=tealfill},           % service / utility
  warn/.style={cls, draw=warnbord, fill=warnfill, font=\sffamily\bfseries\small},  % error state
  ok/.style={cls, draw=okbord, fill=okfill},                  % success state
  neut/.style={cls, draw=neutbord, fill=neutfill},            % neutral
  % ─── size variants ───
  clsSm/.style={cls, minimum height=6mm, minimum width=20mm, font=\sffamily\footnotesize, inner sep=2pt},
  clsLg/.style={cls, minimum height=11mm, minimum width=36mm, font=\sffamily\normalsize},
  absSm/.style={clsSm, draw=absbord, fill=absfill, font=\sffamily\itshape\footnotesize},
  conSm/.style={clsSm, draw=conbord, fill=confill},
  impSm/.style={clsSm, draw=impbord, fill=impfill},
  % ─── circle (existing) ───
  vx/.style={
    circle, draw=absbord, line width=0.4pt,
    minimum size=8mm, fill=absfill, inner sep=0pt
  },
  vxc/.style={vx, draw=conbord, fill=confill},
  vxh/.style={vx, draw=hibord, fill=hifill, font=\sffamily\bfseries},
  % new circle variants
  vxi/.style={vx, draw=impbord, fill=impfill},                % impl green
  vxp/.style={vx, draw=prpbord, fill=prpfill},                % protocol purple
  vxw/.style={vx, draw=warnbord, fill=warnfill},              % warn red
  vxn/.style={vx, draw=nodebord, fill=nodefill},              % neutral
  vxSm/.style={vx, minimum size=6mm, font=\sffamily\footnotesize},
  vxLg/.style={vx, minimum size=11mm, font=\sffamily\normalsize},
  % ─── decision diamond (existing) ───
  q/.style={
    diamond, draw=clibord, line width=0.4pt, fill=clifill,
    inner sep=2pt, aspect=2.4, align=center
  },
  % rounded callout / note (new)
  note/.style={
    rectangle, rounded corners=3pt,
    draw=clibord, line width=0.4pt, fill=clifill,
    inner sep=3pt, font=\sffamily\footnotesize, align=left
  },
  badge/.style={                                              % small inline marker
    rectangle, rounded corners=2pt,
    draw=hibord, line width=0.3pt, fill=hifill,
    inner sep=1.5pt, font=\sffamily\scriptsize\bfseries
  },
  % ─── edges (existing) ───
  inh/.style={                                  % UML inheritance
    -{Triangle[length=3mm,width=3mm,open]},
    draw=edgec, line width=0.5pt
  },
  agg/.style={                                  % UML aggregation
    {Diamond[length=2.5mm,width=2.5mm,open]}-{Stealth[length=2.2mm]},
    draw=edgec, line width=0.45pt
  },
  uses/.style={                                 % directed reference
    -{Stealth[length=2.2mm]}, draw=edgec, line width=0.5pt
  },
  arr/.style={                                  % plain arrow (graphs/trees)
    -{Stealth[length=2mm]}, draw=edgec, line width=0.45pt
  },
  edge/.style={                                 % undirected (trees)
    draw=edgec, line width=0.45pt
  },
  alt/.style={                                  % dashed: similar/alt
    -{Stealth[length=2mm]}, dashed, draw=edgesc, line width=0.45pt
  },
  ret/.style={                                  % dashed creates/returns
    -{Stealth[length=2mm]}, dashed, draw=edgesc, line width=0.45pt
  },
  thr/.style={                                  % threading link (red dashed)
    ->, dashed, draw=conbord, line width=0.5pt
  },
  % ─── new edges ───
  comp/.style={                                 % UML composition (filled diamond)
    {Diamond[length=2.5mm,width=2.5mm]}-{Stealth[length=2.2mm]},
    draw=edgec, line width=0.45pt
  },
  realize/.style={                              % UML realize (impl interface)
    -{Triangle[length=3mm,width=3mm,open]}, dashed,
    draw=edgec, line width=0.5pt
  },
  msgsync/.style={                              % synchronous message (sequence)
    -{Stealth[length=2.6mm]}, draw=edgec, line width=0.55pt
  },
  msgasync/.style={                             % asynchronous message
    -{Stealth[length=2.4mm,inset=0.4mm,fill=edgec]},
    draw=edgec, line width=0.5pt
  },
  msgret/.style={                               % return message (dashed)
    -{Stealth[length=2.2mm]}, dashed, draw=edgesc, line width=0.45pt
  },
  flow/.style={                                 % control flow (thicker, hi-vis)
    -{Stealth[length=2.4mm]}, draw=inkc, line width=0.7pt
  },
  focus/.style={                                % focus / hot path
    -{Stealth[length=2.4mm]}, draw=focusc, line width=0.8pt
  },
  selfloop/.style={                             % self loop on a node
    ->, draw=edgec, line width=0.45pt, in=120, out=60, loop, looseness=8
  },
  % ─── group / region (existing + new) ───
  group/.style={
    rectangle, rounded corners=4pt,
    draw=groupbord, line width=0.4pt,
    inner sep=10pt, fill=black!2
  },
  glabel/.style={font=\sffamily\bfseries\small, anchor=north west, inner sep=2pt},
  panel/.style={                                % stronger panel for layered diagrams
    rectangle, rounded corners=4pt,
    draw=neutbord, line width=0.45pt,
    inner sep=10pt, fill=neutfill!50
  },
  panelTitle/.style={font=\sffamily\bfseries\small, text=inkc, anchor=north west, inner sep=2pt},
  panelSub/.style={font=\sffamily\footnotesize, text=mutec, anchor=north west, inner sep=2pt},
  % ─── typography ───
  caption/.style={font=\sffamily\footnotesize\itshape, text=mutec, align=center},
  legendkey/.style={rectangle, minimum width=10mm, minimum height=3mm, inner sep=0pt, line width=0.3pt},
  % ─── grid / coordinate hints ───
  bggrid/.style={draw=gridc, line width=0.3pt, dotted},
}

% Korean font convenience macro — included by every Hangul diagram
\providecommand{\KOR}{}

\begin{tikzpicture}[blog,
  hdr/.style={font=\sffamily\bfseries\small, text=inkc, align=center},
  bit/.style={rectangle, draw=nodebord, line width=0.4pt, minimum width=5mm, minimum height=8mm, inner sep=1pt, font=\sffamily\scriptsize, align=center, fill=nodefill},
  bitR/.style={bit, fill=confill, draw=conbord},
  bitW/.style={bit, fill=impfill, draw=impbord},
  bitX/.style={bit, fill=clifill, draw=clibord},
  bitN/.style={bit, fill=neutfill, draw=neutbord},
  fldlbl/.style={font=\sffamily\footnotesize, text=inkc, align=center}
]

% Title
\node[hdr] at (0, 5.0) {STM32 USART\_CR1 — 32-bit MMIO register at \texttt{0x40004400 + 0x0C}};

% 32 bits, drawn left-to-right as bit 31 ... bit 0
\foreach \i in {0,...,31} {
  \pgfmathtruncatemacro{\bitn}{31 - \i}
  \node[bit] (b\bitn) at ({-7.75 + \i * 0.5}, 3.0) {\bitn};
}

% Overlay colored field groupings
% Bits 31..29 reserved (neut)
\foreach \i in {31,30,29} {
  \node[bitN] at ({-7.75 + (31-\i) * 0.5}, 3.0) {\i};
}
% Bit 28 M1 — mode bit
\node[bitR] at ({-7.75 + 3 * 0.5}, 3.0) {28};
% Bits 27..15 mix reserved + flags (neut for compactness)
\foreach \i in {27,26,25,24,23,22,21,20,19,18,17,16,15} {
  \node[bitN] at ({-7.75 + (31-\i) * 0.5}, 3.0) {\i};
}
% Bit 14 UE — USART enable (highlight)
\node[bitW] at ({-7.75 + (31-14) * 0.5}, 3.0) {14};
% Bit 13 M0 — mode bit (R group)
\node[bitR] at ({-7.75 + (31-13) * 0.5}, 3.0) {13};
% Bit 12 WAKE
\node[bitN] at ({-7.75 + (31-12) * 0.5}, 3.0) {12};
% Bits 11..10 PCE, PS (parity)
\node[bitX] at ({-7.75 + (31-11) * 0.5}, 3.0) {11};
\node[bitX] at ({-7.75 + (31-10) * 0.5}, 3.0) {10};
% Bit 9 PEIE
\node[bitN] at ({-7.75 + (31-9) * 0.5}, 3.0) {9};
% Bit 8 TXEIE
\node[bitW] at ({-7.75 + (31-8) * 0.5}, 3.0) {8};
% Bit 7 TCIE
\node[bitW] at ({-7.75 + (31-7) * 0.5}, 3.0) {7};
% Bit 6 RXNEIE
\node[bitW] at ({-7.75 + (31-6) * 0.5}, 3.0) {6};
% Bit 5 IDLEIE
\node[bitN] at ({-7.75 + (31-5) * 0.5}, 3.0) {5};
% Bit 4 TE  — transmit enable
\node[bitW] at ({-7.75 + (31-4) * 0.5}, 3.0) {4};
% Bit 3 RE
\node[bitW] at ({-7.75 + (31-3) * 0.5}, 3.0) {3};
% Bit 2 RWU
\node[bitN] at ({-7.75 + (31-2) * 0.5}, 3.0) {2};
% Bit 1 SBK
\node[bitN] at ({-7.75 + (31-1) * 0.5}, 3.0) {1};
% Bit 0 SBRK
\node[bitN] at ({-7.75 + (31-0) * 0.5}, 3.0) {0};

% Field annotations
\draw[edgesc] ({-7.75 + (31-14) * 0.5 + 0.0}, 2.5) -- ({-7.75 + (31-14) * 0.5 + 0.0}, 1.8);
\node[fldlbl] at ({-7.75 + (31-14) * 0.5}, 1.5) {UE};

\draw[edgesc] ({-7.75 + (31-13) * 0.5 + 0.0}, 2.5) -- ({-7.75 + (31-13) * 0.5 + 0.0}, 1.0);
\node[fldlbl] at ({-7.75 + (31-13) * 0.5}, 0.7) {M0};

\draw[edgesc] ({-7.75 + (31-11) * 0.5 + 0.25}, 2.5) -- ({-7.75 + (31-11) * 0.5 + 0.25}, 1.5);
\node[fldlbl] at ({-7.75 + (31-11) * 0.5 + 0.25}, 1.2) {PCE / PS};

\draw[edgesc] ({-7.75 + (31-7) * 0.5 + 0.25}, 2.5) -- ({-7.75 + (31-7) * 0.5 + 0.25}, 1.0);
\node[fldlbl] at ({-7.75 + (31-7) * 0.5 + 0.25}, 0.7) {Interrupt enables};

\draw[edgesc] ({-7.75 + (31-3) * 0.5 + 0.25}, 2.5) -- ({-7.75 + (31-3) * 0.5 + 0.25}, 1.5);
\node[fldlbl] at ({-7.75 + (31-3) * 0.5 + 0.25}, 1.2) {TE / RE};

% Legend
\node[bitR, minimum width=4mm, minimum height=4mm] at (-6.5, -1.2) {};
\node[font=\sffamily\footnotesize, anchor=west, text=inkc] at (-6.2, -1.2) {Mode (M0/M1)};

\node[bitW, minimum width=4mm, minimum height=4mm] at (-2.2, -1.2) {};
\node[font=\sffamily\footnotesize, anchor=west, text=inkc] at (-1.9, -1.2) {Enable bits};

\node[bitX, minimum width=4mm, minimum height=4mm] at (2.0, -1.2) {};
\node[font=\sffamily\footnotesize, anchor=west, text=inkc] at (2.3, -1.2) {Parity};

\node[bitN, minimum width=4mm, minimum height=4mm] at (5.0, -1.2) {};
\node[font=\sffamily\footnotesize, anchor=west, text=inkc] at (5.3, -1.2) {Reserved / 기타};

% Naive vs typed approach text
\node[font=\sffamily\footnotesize, text=conbord, align=left, text width=70mm] at (-4.0, -3.0) {
  \textbf{Naive}: \texttt{USART2->CR1 |= (1 << 13) | (1 << 3);} \\
  $\bullet$ 매직 넘버, 의미 불명
};

\node[font=\sffamily\footnotesize, text=impbord, align=left, text width=70mm] at (3.5, -3.0) {
  \textbf{Typed}: \texttt{cr1.set<UE>().set<TE>();} \\
  $\bullet$ 컴파일러가 검사, 가독성 ↑
};

\end{tikzpicture}
\end{document}
